\documentclass[leqno]{article}
% Shared mark scheme preamble for resources.danieljsmith.org: the pack
% preamble, unchanged, then the scheme's own layer. A scheme's .tex is
% rendered from its YAML (scripts/README.md, Mark schemes), never edited.
\input{\CurrentFilePath/pack_preamble.tex}

% The page: equal side margins, so the table runs the full width of the
% header rule. The footers' right edge follows the text's.
\geometry{left=0.8in, right=0.8in}
\makeatletter
\renewcommand{\djs@footR}[1]{%
  \put(\dimexpr1in+\hoffset+\oddsidemargin+\textwidth\relax,
    -\dimexpr\paperheight-0.32in\relax){\makebox[0pt][r]{#1}}}
\makeatother

\usepackage[export]{adjustbox}
\usepackage{tabularray}
\definecolor{msGrid}{RGB}{176, 172, 164}
\definecolor{msTint}{RGB}{245, 243, 239}

% The header on every page, the front page included, and the booklet label
% under the contents-page title. \djsQuestionPageStyle stays empty.
\newcommand{\djsSchemeHeader}[1]{%
  \djsPackHeader[markscheme]{\textit{Mark scheme} - #1}%
  \renewcommand{\djsPackTitle}{#1}%
  \renewcommand{\djsBookletLabel}{Mark scheme}%
}

% The front page: the contents heading, spaced as tocloft spaces the packs'
% "Questions", then the contents table. A question's name and page number
% both link to its first page.
\newcommand{\msContentsHeading}{%
  \vskip\cftbeforetoctitleskip
  {\noindent\cfttoctitlefont\contentsname\par}%
  \vskip\cftaftertoctitleskip}
\NewTblrEnviron{mscontents}
\SetTblrInner[mscontents]{
  colspec = {Q[wd=26mm, l] X[l] Q[wd=16mm, c] Q[wd=14mm, r]},
  width = \linewidth,
  rowsep = 4pt, colsep = 6pt,
  hlines = {0.4pt, msGrid}, vlines = {0.4pt, msGrid},
  hline{2} = {2.2pt, scarlet},
  row{1} = {font=\small\bfseries},
  row{Z} = {font=\bfseries},
  column{2} = {font=\small},
}
\newcommand{\msContentsQuestion}[1]{\hyperref[q:#1]{Question #1}}
\newcommand{\msContentsPage}[1]{\hyperref[q:#1]{\pageref*{q:#1}}}
\newcommand{\msSep}{\hspace{0.9em}}

% The conventions pages' layout: headings in the documents' own bold, tables
% ruled and padded as the scheme's, roomy lists.
\newcommand{\msHead}[1]{%
  \par\vspace{18pt}%
  {\noindent\bfseries #1\par}%
  \vspace{6pt}}
\NewTblrEnviron{mstblr}
\SetTblrInner[mstblr]{rowsep = 6pt, colsep = 6pt,
  hlines = {0.4pt, msGrid}, vlines = {0.4pt, msGrid}}
\newenvironment{msList}
  {\begin{itemize}[leftmargin=1.2em, itemsep=5pt, topsep=3pt, parsep=0pt]}
  {\end{itemize}}

% A question's first page: the link target and the PDF bookmark, as
% \questionitem gives the packs.
\newcommand{\msQnum}{}
\newcommand{\msQuestion}[1]{%
  \clearpage
  \renewcommand{\msQnum}{#1}%
  \phantomsection\label{q:#1}%
  \addcontentsline{toc}{subsection}{\protect\djsTocEntry{#1}}}

% The scheme's table: Question 16mm, Scheme the rest, Marks 16mm, Notes
% 44mm; every cell ruled in msGrid; the scarlet rule under the head row. A
% long table has no caption or "continued" foot; a continuation page opens
% with a muted line above the repeated head row. The rules across a part's
% rows, the alternative's tint and the total row are the rendered table's.
\NewTblrTheme{ms}{%
  \DefTblrTemplate{firsthead}{default}{}%
  \DefTblrTemplate{middlehead,lasthead}{default}{%
    {\small\itshape\color{djsMuted}Question \msQnum\ continued\par}%
    \vspace{6pt}}%
  \DefTblrTemplate{firstfoot,middlefoot,lastfoot}{default}{}%
}
\NewTblrEnviron{msscheme}
\SetTblrOuter[msscheme]{long, theme=ms, entry=none, label=none}
\SetTblrInner[msscheme]{
  colspec = {Q[wd=16mm, l] X[l] Q[wd=16mm, c] Q[wd=44mm, l]},
  width = \linewidth,
  rowsep = 6pt, colsep = 6pt,
  rowhead = 1,
  vlines = {0.4pt, msGrid},
  hline{1} = {0.4pt, msGrid}, hline{2} = {2.2pt, scarlet},
  hline{Z} = {0.4pt, msGrid},
  column{1} = {font=\bfseries}, column{3} = {font=\bfseries},
  column{4} = {font=\small},
  row{1} = {rowsep=4pt, font=\small\bfseries},
}

% Two marks on a row: each note starts with its code.
\newcommand{\msNote}[2]{\textbf{#1}\enspace #2}
% A full-row figure: the worked solution's own picture, scaled uniformly and
% only down, to the table less the Question column. The picture is drawn
% with the line width it has in the solutions document, at its list depth
% there (the optional argument, 1 by default: a part's working), so one
% sized by \linewidth keeps its size; tabularray sets the cell's \linewidth
% to the cell's measured width, which would shrink it at every pass. The
% widths are measured once, in a solution box at the packs' text width
% (pack_preamble.tex's margins, 0.8in and 1in).
\newlength{\msSpanW}
\makeatletter
\newcommand{\ms@keepLW}[1]{%
  \expandafter\xdef\csname ms@solnLW@#1\endcsname{\the\linewidth}}
\AtBeginDocument{%
  \setlength{\msSpanW}{\dimexpr\textwidth-16mm-28pt\relax}%
  \setbox\z@\vbox{%
    \begin{minipage}{\dimexpr\paperwidth-1.8in\relax}
      \begin{djsSolution}
        \ms@keepLW{0}%
        \begin{enumerate}\item \ms@keepLW{1}%
        \begin{enumerate}\item \ms@keepLW{2}%
        \begin{enumerate}\item \ms@keepLW{3}%
        \end{enumerate}\end{enumerate}\end{enumerate}
      \end{djsSolution}
    \end{minipage}}}
\newcommand{\msFigure}[2][1]{%
  \adjustbox{max width=\msSpanW, max totalheight=0.6\textheight, valign=t}{%
    \hsize=\csname ms@solnLW@#1\endcsname
    \linewidth=\hsize
    #2}}
\makeatother
% A figure mark's feature, in the Scheme column.
\newcommand{\msFeature}[1]{\textit{#1}}
% An alternative's title row.
\newcommand{\msAlternative}[1]{\textit{Alternative:} #1}
% Drafts only: a queried mark's code, the query's why under its note, and a
% flag in its part's place.
\newcommand{\msQueried}[2]{#1~\(\to\) #2?}
\newcommand{\msQuery}[1]{{\itshape\color{djsMuted}Query: #1}}
\newcommand{\msFlag}[2]{\textbf{#1:}\enspace #2}


\djsSchemeHeader{Improper Integrals and Mean Values}
\djsSitePath{Core Pure}{Further Calculus}

\djsLastUpdated{26 October 2026}
\begin{document}
\djsSetQuestionCount{16}
\djsAddMarks{1}{4}
\djsAddMarks{2}{5}
\djsAddMarks{3}{8}
\djsAddMarks{4}{6}
\djsAddMarks{5}{6}
\djsAddMarks{6}{8}
\djsAddMarks{7}{7}
\djsAddMarks{8}{7}
\djsAddMarks{9}{7}
\djsAddMarks{10}{8}
\djsAddMarks{11}{7}
\djsAddMarks{12}{6}
\djsAddMarks{13}{9}
\djsAddMarks{14}{9}
\djsAddMarks{15}{8}
\djsAddMarks{16}{11}
\vspace*{6pt}
\djsContentsTitle
\vspace{30pt}
\msContentsHeading
\begin{mscontents}{}
Question & Marks by part & Marks & Page \\
\msContentsQuestion{1} & & 4 & \msContentsPage{1} \\
\msContentsQuestion{2} & & 5 & \msContentsPage{2} \\
\msContentsQuestion{3} & & 8 & \msContentsPage{3} \\
\msContentsQuestion{4} & & 6 & \msContentsPage{4} \\
\msContentsQuestion{5} & & 6 & \msContentsPage{5} \\
\msContentsQuestion{6} & (a)\,1\msSep (b)\,3\msSep (c)\,4 & 8 & \msContentsPage{6} \\
\msContentsQuestion{7} & (a)\,2\msSep (b)\,2\msSep (c)\,3 & 7 & \msContentsPage{7} \\
\msContentsQuestion{8} & (a)\,4\msSep (b)\,3 & 7 & \msContentsPage{8} \\
\msContentsQuestion{9} & & 7 & \msContentsPage{9} \\
\msContentsQuestion{10} & & 8 & \msContentsPage{10} \\
\msContentsQuestion{11} & (a)\,1\msSep (b)\,6 & 7 & \msContentsPage{11} \\
\msContentsQuestion{12} & (a)\,1\msSep (b)\,5 & 6 & \msContentsPage{12} \\
\msContentsQuestion{13} & (a)\,1\msSep (b)\,4\msSep (c)\,4 & 9 & \msContentsPage{13} \\
\msContentsQuestion{14} & (a)\,6\msSep (b)\,1\msSep (c)\,2 & 9 & \msContentsPage{14} \\
\msContentsQuestion{15} & & 8 & \msContentsPage{15} \\
\msContentsQuestion{16} & (a)\,1\msSep (b)(i)\,3\msSep (b)(ii)\,6\msSep (b)(iii)\,1 & 11 & \msContentsPage{16} \\
\SetCell[c=2]{r} Total & & 116 & \\
\end{mscontents}
\newpage
% The conventions pages, shared by every pack's mark scheme: one source of
% truth, \input after the front page, with no strand or pack logic. The keys
% on the first side, the rules on the second. The scheme layer defines
% \msHead, msList and mstblr.
\vspace*{6pt}%
{\noindent\cfttoctitlefont Reading the mark scheme\par}%
\vspace{10pt}%
These rules apply throughout, unless a note says otherwise.

\msHead{Marks}
Every mark is worth one, so \textbf{M1 A1} on a row is two marks. Only marks
coded ft follow through.\par
\vspace{10pt}%
\begin{mstblr}{
  colspec = {Q[wd=15mm, l] X[l]},
  width = \linewidth,
  hline{2} = {2.2pt, scarlet},
  column{1} = {font=\bfseries},
  row{1} = {font=\small\bfseries, rowsep=4pt},
}
Code & Meaning \\
M1 & Method: a correct method, applied to this question. A slip doesn't cost it if the method is clear. A formula must be used, not just quoted. \\
dM1 & Dependent method: as M1, but only if the method mark before it was earned, unless its note names another. \\
A1 & Accuracy: a correct answer or step. It needs the method mark it follows. \\
A1ft & As A1, following through from an earlier answer that the note names. \\
A1~AG & The last mark of a part whose answer is printed in the question: complete working, with no errors. \\
B1 & A correct value, fact or statement. It needs no method mark. \\
B1ft & As B1, following through from an earlier answer that the note names. \\
E1 & A written explanation, interpretation or conclusion. Concluding a proof by induction needs every earlier mark. \\
\end{mstblr}

\msHead{Abbreviations}
ft and AG qualify a code in the Marks column. The others appear in the
notes.\par
\vspace{10pt}%
\begin{mstblr}{
  colspec = {Q[wd=17mm, l] X[l] Q[wd=15mm, l] X[l]},
  width = \linewidth,
  column{1} = {font=\bfseries}, column{3} = {font=\bfseries},
}
ft & follow through & AG & answer given in the question \\
cao & correct answer only, with no follow-through & cso & correct solution only: no errors in the part \\
oe & or an equivalent form & awrt & any answer that rounds to this value \\
soi & seen, or implied by later working & isw & ignore working after a correct answer \\
SC & special case: credit outside the main method & their & the student's own earlier value \\
condone & a slip that doesn't cost the mark & [\;\,] & working the mark doesn't need \\
s.f. & significant figures & d.p. & decimal places \\
\end{mstblr}

\newpage
\msHead{Reading the scheme}
\begin{msList}
\item Each row is a line of working, with the marks it earns beside it.
  Where a row carries two marks, each note starts with its code.
\item \textbf{[\textit{n}]} ends each part with its total.
\item A shaded \textit{Alternative} is a second method for the same total.
  Mark an answer by the method it follows.
\item Any other valid method earns the equivalent marks, unless the question
  names a method or says \textit{hence}.
\item A sketch earns one B1 for each feature named below its figure.
\end{msList}

\msHead{Accuracy and answers}
\begin{msList}
\item Unless a question says otherwise, give answers to 3 significant
  figures, and angles in degrees to 1 decimal place. Angles in radians take
  3 significant figures.
\item A final answer with more figures is accepted if it rounds to the
  scheme's value; one with fewer is not. An answer spoiled by rounding too
  early loses its accuracy mark.
\item Where a question asks for an exact answer, a decimal loses the final
  accuracy mark. Otherwise either is accepted.
\item Wrong or missing units cost nothing unless the question asks for them.
\item A correct answer with no working earns full marks, unless the question
  needs the working: show that, prove, hence, a method it names, or an exact
  answer. A wrong answer with no working earns nothing. Calculator answers
  are treated the same way.
\item A correct answer, once reached, isn't spoiled by later working (isw).
\end{msList}

\msHead{Value of $\boldsymbol{g}$}
\begin{msList}
\item Take $g = 9.8\,\mathrm{m\,s^{-2}}$ unless a question states a value,
  which is then the only one accepted. Answers that use $g = 9.8$ may be
  given to 2 or 3 significant figures.
\item Where a question doesn't state $g$, an answer worked with $g = 10$
  throughout earns the same marks, keeping any value printed in the
  question. The worked solutions list the final answers with $g = 10$.
\end{msList}


\msQuestion{1}
\begin{msscheme}{
  hline{3-4} = {2-Z}{0.4pt, msGrid},
  hline{5} = {0.4pt, msGrid},
  row{5} = {font=\bfseries},
}
Question & Scheme & Marks & Notes \\
1 & {\(\int \frac{x}{x^2+1}\,\mathrm{d}x = \tfrac12\ln(x^2+1)\)} & M1 A1 & {\msNote{M1}{Integrates to the form \(k\ln(x^2+1)\), or \(k\ln u\) after substituting \(u = x^2+1\).}\\ \msNote{A1}{oe, such as \(\frac12\ln u\). Condone modulus signs.}} \\
& {\(\frac{1}{3-1}\left(\their{\tfrac12}\ln 10 - \their{\tfrac12}\ln 2\right)\)} & dM1 & {Substitutes limits 1 and 3 (or 2 and 10 in \(u\)) into their \(k\ln(x^2+1)\), subtracts, and divides by 2.} \\
& {\(\tfrac14\ln 5\)} & A1 & {Must be in the form \(p\ln q\); accept \(\frac18\ln 25\). A0 for \(\frac14(\ln 10 - \ln 2)\) or a decimal.} \\
\nopagebreak \SetCell[c=2]{r} Total for Question 1 & & 4 & \\
\end{msscheme}

\msQuestion{2}
\begin{msscheme}{
  hline{3-6} = {2-Z}{0.4pt, msGrid},
  hline{7} = {0.4pt, msGrid},
  row{7} = {font=\bfseries},
}
Question & Scheme & Marks & Notes \\
2 & {\(\int x(x^2+5)^{-\frac32}\,\mathrm{d}x = k(x^2+5)^{-\frac12}\)} & M1 & {Integrates to the form \(k(x^2+5)^{-\frac12}\), \(k \neq 0\), by reverse chain rule or substitution \(u = x^2+5\).} \\
& {\(\phantom{\int} = -(x^2+5)^{-\frac12}\ (+c)\)} & A1 & {oe, such as \(-\frac{1}{\sqrt{u}}\) with \(u = x^2+5\).} \\
& {\(\int_2^t \ldots\,\mathrm{d}x = \left[-\frac{1}{\sqrt{x^2+5}}\right]_2^t\)\\[4pt] \(\phantom{\int_2^t} = -\frac{1}{\sqrt{t^2+5}} + \frac13\)} & M1 & {Replaces \(\infty\) by \(t\) (any letter) and substitutes \(t\) and 2, or corresponding limits in \(u\), into their \(k(x^2+5)^{-\frac12}\), subtracting.} \\
& {As \(t\to\infty\), \(\frac{1}{\sqrt{t^2+5}}\to 0\)} & dM1 & {States that their term in \(t\) tends to 0 as \(t\to\infty\), in limit notation or words. Not for substituting \(\infty\) directly.} \\
& {\(\int_2^\infty \ldots\,\mathrm{d}x = \frac13 - 0 = \frac13\)} & A1~AG & {cso: correct antiderivative and limits leading without error to \(\frac13\), with the limiting process shown.} \\
\nopagebreak \SetCell[c=2]{r} Total for Question 2 & & 5 & \\
\end{msscheme}

\msQuestion{3}
\begin{msscheme}{
  hline{3-8} = {2-Z}{0.4pt, msGrid},
  hline{9} = {0.4pt, msGrid},
  row{9} = {font=\bfseries},
}
Question & Scheme & Marks & Notes \\
3 & {\(\dfrac{4}{x(x+1)(x+2)}=\dfrac{A}{x}+\dfrac{B}{x+1}+\dfrac{C}{x+2}\)} & M1 & {Uses the correct form with three constants. Not for a form missing a term.} \\
& {\(4=A(x+1)(x+2)+Bx(x+2)+Cx(x+1)\)\\[4pt] \(\Rightarrow A=\ldots,\ B=\ldots,\ C=\ldots\)} & dM1 & {Finds all three constants, by substituting values or equating coefficients. Condone arithmetic slips.} \\
& {\(\dfrac{2}{x}-\dfrac{4}{x+1}+\dfrac{2}{x+2}\)} & A1 & \\
& {\(2\ln x-4\ln(x+1)+2\ln(x+2)\)} & M1 A1 & {\msNote{M1}{Integrates at least two of their partial fractions of the form \(\frac{a}{x+b}\) to \(a\ln(x+b)\). Modulus signs not needed.}\\ \msNote{A1}{All three correct; need not be combined.}} \\
& {\(\left[2\ln\dfrac{x(x+2)}{(x+1)^2}\right]_1^{t}\)\\[4pt] \(\phantom{A}=2\ln\dfrac{t(t+2)}{(t+1)^2}-2\ln\dfrac{3}{4}\)} & M1 & {Replaces \(\infty\) by \(t\) or uses a limit, combines their log terms into one logarithm of a function of \(t\), and substitutes both limits. Condone sign slips.} \\
& {As \(t\to\infty\), \(\dfrac{t^2+2t}{t^2+2t+1}=\dfrac{1+\frac{2}{t}}{1+\frac{2}{t}+\frac{1}{t^2}}\to1\)\\[4pt] \(\Rightarrow 2\ln(\ldots)\to0\), so the integral converges} & dM1 & {Shows the argument of the combined logarithm tends to 1 (e.g. dividing by \(t^2\)), so its log tends to 0, and concludes convergence. Not for 0 written without a limit argument.} \\
& {\(\ln\dfrac{16}{9}\)} & A1 & {cso. Must be in the form \(\ln k\), or \(k=\frac{16}{9}\) stated; \(2\ln\frac43\) alone does not score.} \\
\nopagebreak \SetCell[c=2]{r} Total for Question 3 & & 8 & \\
\end{msscheme}

\msQuestion{4}
\begin{msscheme}{
  hline{3-6} = {2-Z}{0.4pt, msGrid},
  hline{7} = {0.4pt, msGrid},
  row{7} = {font=\bfseries},
}
Question & Scheme & Marks & Notes \\
4 & {\(y = \dfrac{1+\sin x}{1-\sin^2 x} = \dfrac{1+\sin x}{\cos^2 x}\)} & M1 & {Multiplies numerator and denominator by \(1+\sin x\) and uses \(1-\sin^2 x=\cos^2 x\) to reach a denominator of \(\cos^2 x\).} \\
& {\(y = \sec^2 x + \sec x\tan x\)} & A1 & \\
& {\(\displaystyle\int \left(\sec^2 x + \sec x\tan x\right)\mathrm{d}x = \tan x + \sec x\)} & M1 A1 & {\msNote{M1}{Integrates their expression of the form \(p\sec^2 x + q\sec x\tan x\) to the form \(a\tan x + b\sec x\).}\\ \msNote{A1}{A1: Correct \(\tan x + \sec x\).}} \\
& {\(\dfrac{6}{\pi}\left[\left(\tfrac{1}{\sqrt3}+\tfrac{2}{\sqrt3}\right)-(0+1)\right]\)} & dM1 & {Substitutes both limits into their integral, subtracts, and divides by \(\frac{\pi}{6}\). Need not be simplified.} \\
& {\(\dfrac{6\left(\sqrt3-1\right)}{\pi}\)} & A1 & {cso. Any equivalent exact form.} \\
\nopagebreak \SetCell[c=2]{r} Total for Question 4 & & 6 & \\
\end{msscheme}

\msQuestion{5}
\begin{msscheme}{
  hline{3-6} = {2-Z}{0.4pt, msGrid},
  hline{7} = {0.4pt, msGrid},
  row{7} = {font=\bfseries},
}
Question & Scheme & Marks & Notes \\
5 & {\(\frac{1}{2}\int_0^2 40\tanh^2\left(\tfrac12 t\right)\mathrm{d}t\)\\[4pt] \(\phantom{A} = 20\int_0^2 \left(1-\operatorname{sech}^2\left(\tfrac12 t\right)\right)\mathrm{d}t\)} & M1 & {Uses a hyperbolic identity, such as \(\tanh^2 x = 1-\operatorname{sech}^2 x\), to replace \(\tanh^2\left(\frac12 t\right)\) in the integral of \(p\). Condone a sign slip in the identity.} \\
& {\(20\left[t-2\tanh\left(\tfrac12 t\right)\right]_0^2\)} & M1 A1 & {\msNote{M1}{Integrates to the form \(\alpha t+\beta\tanh\left(\frac12 t\right)\).}\\ \msNote{A1}{Correct, with factor 20, or 40 if not yet divided by 2.}} \\
& {\(20\left(2-2\tanh 1\right)-20(0) \ [=40(1-\tanh 1)]\)} & dM1 & {Substitutes both limits into their integral and divides by 2. Lower limit giving 0 need not be shown.} \\
& {\(\tanh 1=\frac{e-e^{-1}}{e+e^{-1}} \ \left[=\frac{e^2-1}{e^2+1}\right]\)} & M1 & {Writes \(\tanh 1\), or \(\sinh 1\) and \(\cosh 1\), in exponential form in their expression. Need not be simplified.} \\
& {\(\frac{80}{e^2+1}\)} & A1 & {cso. Only this form scores, not e.g. \(\frac{80e^{-2}}{1+e^{-2}}\) or \(40(1-\tanh 1)\).} \\
\nopagebreak \SetCell[c=2]{r} Total for Question 5 & & 6 & \\
\end{msscheme}

\msQuestion{6}
\begin{msscheme}{
  hline{3,5-6,8-11} = {2-Z}{0.4pt, msGrid},
  hline{4,7,12} = {0.4pt, msGrid},
  row{12} = {font=\bfseries},
}
Question & Scheme & Marks & Notes \\
6(a) & {Upper limit is infinite, so the interval of integration is unbounded} & E1 & {States that the upper limit is infinite, or the interval unbounded. Not for a reason about the integrand alone.} \\
\nopagebreak & & [1] & \\
6(b) & {\(u = \ln x,\ \frac{\mathrm{d}v}{\mathrm{d}x} = x^{-3} \Rightarrow v = -\frac{1}{2x^2}\)\\[4pt] \(\int \frac{\ln x}{x^3}\,\mathrm{d}x = -\frac{\ln x}{2x^2} + \int \frac{1}{2x^3}\,\mathrm{d}x\)} & M1 A1 & {\msNote{M1}{Integrates by parts with \(u = \ln x\), reaching the form \(px^{-2}\ln x \pm \int qx^{-3}\,\mathrm{d}x\). Condone sign and coefficient slips.}\\ \msNote{A1}{Correct first stage; need not be simplified.}} \\
& {\(\phantom{\int \frac{\ln x}{x^3}\,\mathrm{d}x} = -\frac{\ln x}{2x^2} - \frac{1}{4x^2} + c\)} & A1 & {oe. Needs \(+c\).} \\
\nopagebreak & & [3] & \\
6(c) & {\(\int_e^t \frac{\ln x}{x^3}\,\mathrm{d}x = \left[\their{-\frac{\ln x}{2x^2} - \frac{1}{4x^2}}\right]_e^t\)} & M1 & {Replaces \(\infty\) by \(t\) and substitutes \(t\) and \(e\) into their (b) of the form \(px^{-2}\ln x + qx^{-2}\). Not for substituting \(\infty\) directly.} \\
& {\(\phantom{\int_e^t} = \frac{3}{4e^2} - \frac{\ln t}{2t^2} - \frac{1}{4t^2}\)} & A1 & {oe; \(\ln e\) need not be evaluated.} \\
& {As \(t\to\infty\), \(\frac{\ln t}{t^2}\to 0\) and \(\frac{1}{t^2}\to 0\)} & dM1 & {Lets \(t\to\infty\), stating explicitly that \(\frac{\ln t}{t^2}\to 0\) (oe) and \(\frac{1}{t^2}\to 0\). Not for replacing the terms by 0 without a statement.} \\
& {Converges: \(\int_e^\infty \frac{\ln x}{x^3}\,\mathrm{d}x = \frac{3}{4e^2}\)} & A1 & {cso. Exact value oe, with a statement that the integral converges or the limit is finite.} \\
\nopagebreak & & [4] & \\
\nopagebreak \SetCell[c=2]{r} Total for Question 6 & & 8 & \\
\end{msscheme}

\msQuestion{7}
\begin{msscheme}{
  hline{3-4,6-7,9-11} = {2-Z}{0.4pt, msGrid},
  hline{5,8,12} = {0.4pt, msGrid},
  row{12} = {font=\bfseries},
}
Question & Scheme & Marks & Notes \\
7(a) & {\(9(x^2-2x)+25 = 9\left[(x-1)^2-1\right]+25\)} & M1 & {Takes out 9 from the \(x\) terms and completes the square to \(9(x\pm1)^2+q\), or compares coefficients of \(a(x+p)^2+q\). Condone a slip in \(q\).} \\
& {\(= 9(x-1)^2+16\)} & A1 & {Or \(a=9\), \(p=-1\), \(q=16\) stated.} \\
\nopagebreak & & [2] & \\
7(b) & {\(\int\frac{\mathrm{d}x}{\sqrt{9(x-1)^2+16}}\)\\[4pt] \(\phantom{=} = \frac13\int\frac{\mathrm{d}x}{\sqrt{(x-1)^2+\left(\frac43\right)^2}}\)} & M1 & {Uses their (a), of the form \(9(x+p)^2+q\) with \(q>0\), with the arsinh result to reach \(\lambda\operatorname{arsinh}(\mu(x+p))\) or a logarithmic equivalent. Condone slips in \(\lambda\), \(\mu\).} \\
& {\(= \frac13\operatorname{arsinh}\left(\frac{3(x-1)}{4}\right)+c\)} & A1 & {oe, including a logarithmic form. Needs \(+c\).} \\
\nopagebreak & & [2] & \\
7(c) & {\(\frac12\times\frac13\left[\operatorname{arsinh}\left(\frac{3(x-1)}{4}\right)\right]_0^2\)\\[4pt] \(\phantom{=} = \frac16\left(\operatorname{arsinh}\frac34-\operatorname{arsinh}\left(-\frac34\right)\right)\)} & M1 & {Substitutes 0 and 2 into their (b), of the form \(\lambda\operatorname{arsinh}(\mu(x+p))\) or a logarithmic equivalent, subtracts and divides by 2. Condone sign slips.} \\
& {\(\operatorname{arsinh}\frac34 = \ln\left(\frac34+\sqrt{\frac9{16}+1}\right) = \ldots\)} & M1 & {Uses \(\operatorname{arsinh}t=\ln\left(t+\sqrt{t^2+1}\right)\) on at least one of their values, or evaluates a logarithmic form from (b) at a limit.} \\
& {\(\frac16(\ln 2+\ln 2) = \frac13\ln 2\)} & A1 & {cso. Any form \(A\ln B\) with \(B\) an integer, such as \(\frac16\ln4\). Not \(\frac16\left(\ln2-\ln\frac12\right)\).} \\
\nopagebreak & & [3] & \\
\nopagebreak \SetCell[c=2]{r} Total for Question 7 & & 7 & \\
\end{msscheme}

\msQuestion{8}
\begin{msscheme}{
  hline{3-6,8-10} = {2-Z}{0.4pt, msGrid},
  hline{7,11} = {0.4pt, msGrid},
  row{11} = {font=\bfseries},
}
Question & Scheme & Marks & Notes \\
8(a) & {\(\int \ln(1+x)\,\mathrm{d}x = x\ln(1+x) - \int \frac{x}{1+x}\,\mathrm{d}x\)} & M1 & {Applies parts with \(u = \ln(1+x)\), reaching \(px\ln(1+x) - \int \frac{qx}{1+x}\,\mathrm{d}x\), or the equivalent after substituting \(t = 1+x\). Condone sign slips.} \\
& {\(\phantom{\int} = x\ln(1+x) - x + \ln(1+x)\)} & A1 & {oe, such as \((1+x)\ln(1+x) - x\) or \(t\ln t - t\).} \\
& {\(\frac{1}{e-1}\Big[(1+x)\ln(1+x) - x\Big]_0^{e-1}\)\\[4pt] \(\phantom{A} = \frac{1}{e-1}\big((e\ln e - (e-1)) - 1\times\ln 1\big)\)} & dM1 & {Substitutes both limits, or 1 and \(e\) after substituting \(t = 1+x\), into their antiderivative and divides by \(e-1\).} \\
& {\(\phantom{A} = \frac{1}{e-1}\)} & A1~AG & {Correct substituted expression, using \(\ln e = 1\) and \(\ln 1 = 0\), leading without error to the printed result.} \\
\nopagebreak & & [4] & \\
8(b) & {\(x = e^{y} - 1\)} & B1 & \\
& {\(\frac{1}{1-0}\int_0^1 (e^{y} - 1)\,\mathrm{d}y = \Big[e^{y} - y\Big]_0^1 = \ldots\)} & M1 & {Integrates their \(x\) in terms of \(y\), at least one term correct, and substitutes \(y = 1\) and \(y = 0\). Division by \(1-0\) need not be shown.} \\
& {\(e - 2\)} & A1 & \\
\nopagebreak & & [3] & \\
\nopagebreak \SetCell[c=2]{r} Total for Question 8 & & 7 & \\
\end{msscheme}

\msQuestion{9}
\begin{msscheme}{
  hline{3-8} = {2-Z}{0.4pt, msGrid},
  hline{9} = {0.4pt, msGrid},
  row{9} = {font=\bfseries},
}
Question & Scheme & Marks & Notes \\
9 & {\(\frac{n}{2\pi}\int_0^{2\pi/n}\left(a^2\sin^2 nt + 2ab\sin nt\sin 2nt\right.\)\\[4pt] \(\phantom{\frac{n}{2\pi}\int_0^{2\pi/n}}\left.{}+ b^2\sin^2 2nt\right)\mathrm{d}t\)} & M1 & {Forms the integral of \(I^2\), expanded into three terms, from \(0\) to \(\frac{2\pi}{n}\), with or without \(\frac{n}{2\pi}\). Condone a slip in the cross term's coefficient.} \\
& {\(\int\sin^2 nt\,\mathrm{d}t = \int\tfrac12(1-\cos 2nt)\,\mathrm{d}t\)\\[4pt] \(\phantom{\int\sin^2 nt\,\mathrm{d}t} = \frac{t}{2} - \frac{\sin 2nt}{4n}\)} & M1 & {Uses \(\sin^2 x = \frac12(1\pm\cos 2x)\) for at least one square and integrates to the form \(pt + q\sin knt\). Condone sign and coefficient slips.} \\
& {\(\int_0^{2\pi/n}\sin^2 nt\,\mathrm{d}t = \int_0^{2\pi/n}\sin^2 2nt\,\mathrm{d}t = \frac{\pi}{n}\)} & A1 & \\
& {\(\int\sin nt\sin 2nt\,\mathrm{d}t = \int 2\sin^2 nt\cos nt\,\mathrm{d}t\)\\[4pt] \(\phantom{\int\sin nt\sin 2nt\,\mathrm{d}t} = \frac{2\sin^3 nt}{3n}\)} & M1 & {Integrates the cross term by a valid method, e.g. to the form \(k\sin^3 nt\), or by product-to-sum to \(p\sin nt + q\sin 3nt\). Condone coefficient slips.} \\
& {\(\left[\frac{2\sin^3 nt}{3n}\right]_0^{2\pi/n} = 0\)} & A1 & {Limits substituted into a correct integral; a bare statement that the integral is zero does not score.} \\
& {\(\frac{n}{2\pi}\left(a^2\times\their{\tfrac{\pi}{n}} + 2ab\times\their{0} + b^2\times\their{\tfrac{\pi}{n}}\right)\)\\[4pt] \(\Rightarrow \text{RMS} = \sqrt{\ldots}\)} & dM1 & {Combines their three integrals with their coefficients, multiplies by \(\frac{n}{2\pi}\) and takes the square root. Depends on the first M1.} \\
& {\(\frac{a^2+b^2}{2} \Rightarrow \text{RMS} = \sqrt{\frac{a^2+b^2}{2}}\)} & A1~AG & {cso: correct integrals combined without error to the printed result.} \\
\nopagebreak \SetCell[c=2]{r} Total for Question 9 & & 7 & \\
\end{msscheme}

\msQuestion{10}
\begin{msscheme}{
  hline{3-8} = {2-Z}{0.4pt, msGrid},
  hline{9} = {0.4pt, msGrid},
  row{9} = {font=\bfseries},
}
Question & Scheme & Marks & Notes \\
10 & {\(\dfrac{x+3}{x^{2}(x+1)}=\dfrac{A}{x}+\dfrac{B}{x^2}+\dfrac{C}{x+1}\)} & M1 & {Uses partial fractions with separate terms in \(x\) and \(x^2\) for the repeated factor, and a term in \(x+1\).} \\
& {\(x+3=Ax(x+1)+B(x+1)+Cx^2\)\\[4pt] \(\Rightarrow A=\ldots,\ B=\ldots,\ C=\ldots\)} & dM1 & {Clears fractions and finds all three constants, by substitution or equating coefficients. Condone arithmetic slips.} \\
& {\(-\dfrac{2}{x}+\dfrac{3}{x^2}+\dfrac{2}{x+1}\)} & A1 & \\
& {\(\left[-2\ln x-\dfrac{3}{x}+2\ln(x+1)\right]\)} & M1 A1 & {\msNote{M1}{Integrates their partial fractions to the form \(p\ln x+\frac{q}{x}+r\ln(x+1)\). Condone sign and coefficient slips.}\\ \msNote{A1}{Correct antiderivative; modulus signs optional.}} \\
& {\(\lim_{t\to\infty}\left[2\ln\left(\dfrac{x+1}{x}\right)-\dfrac{3}{x}\right]_{1}^{t}\)\\[4pt] \(=\lim_{t\to\infty}\left(2\ln\left(1+\tfrac{1}{t}\right)-\tfrac{3}{t}-2\ln 2+3\right)\)} & M1 & {Uses an upper limit \(t\), combines their two log terms into one logarithm, before or after substituting, and substitutes both limits. Not if \(\infty\) is substituted.} \\
& {\(t\to\infty:\ \ln\left(1+\tfrac{1}{t}\right)\to\ln 1=0,\ \tfrac{3}{t}\to 0\)} & dM1 & {States, as \(t\to\infty\), that their combined log term tends to \(\ln 1=0\) and their \(\frac{k}{t}\) term tends to 0.} \\
& {\(\displaystyle\int_{1}^{\infty}\dfrac{x+3}{x^{2}(x+1)}\,\mathrm{d}x=3-2\ln 2=3-\ln 4\)} & A1~AG & {cso, from correct limits leading without error to \(3-\ln 4\). Withhold if \(\infty\) is substituted or \(\ln\infty\) terms are cancelled anywhere.} \\
\nopagebreak \SetCell[c=2]{r} Total for Question 10 & & 8 & \\
\end{msscheme}

\msQuestion{11}
\begin{msscheme}{
  hline{3,5-10} = {2-Z}{0.4pt, msGrid},
  hline{4,11} = {0.4pt, msGrid},
  row{11} = {font=\bfseries},
}
Question & Scheme & Marks & Notes \\
11(a) & {Upper limit is infinite, so the interval of integration is unbounded} & E1 & {States that the upper limit is infinite, or the interval is unbounded. Not for a reference to the integrand alone.} \\
\nopagebreak & & [1] & \\
11(b) & {\(\frac{10}{(x+2)(3x+1)} = \frac{A}{x+2}+\frac{B}{3x+1}\)\\[4pt] \(10 = A(3x+1)+B(x+2) \Rightarrow A = \ldots,\ B = \ldots\)} & M1 & {Uses the correct partial-fraction form and a method for both constants, by substitution or comparing coefficients.} \\
& {\(A=-2,\ B=6\)} & A1 & \\
& {\(\int\left(-\frac{2}{x+2}+\frac{6}{3x+1}\right)\mathrm{d}x\)\\[4pt] \(\phantom{A} = p\ln(x+2)+q\ln(3x+1)\)} & M1 & {Integrates their two partial fractions to the form \(p\ln(x+2)+q\ln(3x+1)\). Condone coefficient slips and missing modulus signs.} \\
& {\(= 2\ln\left(\frac{3x+1}{x+2}\right)\)} & A1 & {Correct single logarithm, oe; may first be seen after the limits are substituted.} \\
& {\(\left[2\ln\frac{3x+1}{x+2}\right]_0^t = 2\ln\frac{3t+1}{t+2} - 2\ln\frac12\)\\[4pt] \(\frac{3t+1}{t+2} = \frac{3+\frac1t}{1+\frac2t} \to \ldots\) as \(t\to\infty\)} & dM1 & {Replaces \(\infty\) by a variable, substitutes both limits and finds the limit of their logarithm's argument as it tends to infinity. Not for substituting \(\infty\) directly.} \\
& {\(2\ln 3 + 2\ln 2 = 2\ln 6\)} & A1 & {cso. Any correct \(a\ln b\) with integer \(a\) and rational \(b\), or \(a=2\), \(b=6\) stated.} \\
\nopagebreak & & [6] & \\
\nopagebreak \SetCell[c=2]{r} Total for Question 11 & & 7 & \\
\end{msscheme}

\msQuestion{12}
\begin{msscheme}{
  hline{3,5-9} = {2-Z}{0.4pt, msGrid},
  hline{4,10} = {0.4pt, msGrid},
  row{10} = {font=\bfseries},
}
Question & Scheme & Marks & Notes \\
12(a) & {At \(x=1\), \(\sqrt[3]{x^2-1}=0\): integrand undefined\\[4pt] \(x=1\) lies in \([0,3]\)} & E1 & {States the integrand is undefined or unbounded (or denominator zero) at \(x=1\), and that 1 lies within the interval of integration.} \\
\nopagebreak & & [1] & \\
12(b) & {\(\int x\left(x^2-1\right)^{-\frac13}\,\mathrm{d}x = k\left(x^2-1\right)^{\frac23}\)} & M1 & {Integrates to the form \(k\left(x^2-1\right)^{\frac23}\), or \(ku^{\frac23}\) with \(u=x^2-1\), by inspection or substitution.} \\
& {\(\phantom{\int} = \frac34\left(x^2-1\right)^{\frac23}\ (+c)\)} & A1 & {oe, including \(\frac34u^{\frac23}\) with \(u=x^2-1\).} \\
& {Splits at \(x=1\):\\[4pt] \(\lim_{t\to1^-}\left[\frac34\left(x^2-1\right)^{\frac23}\right]_0^t\)\\[4pt] \(\phantom{=} + \lim_{s\to1^+}\left[\frac34\left(x^2-1\right)^{\frac23}\right]_s^3\)} & M1 & {Splits at \(x=1\) into two integrals, each written with a limit of a variable endpoint tending to 1 from the appropriate side.} \\
& {\(= \left(0-\frac34\right) + \left(3-0\right)\)} & dM1 & {Substitutes the endpoints into their antiderivative in both integrals and evaluates the limits, with the terms at the variable endpoint tending to 0.} \\
& {Both limits finite, so converges; \(=\frac94\)} & A1 & {cso. \(\frac94\) oe with a statement that the integral converges or both limits are finite. A0 without the limit argument.} \\
\nopagebreak & & [5] & \\
\nopagebreak \SetCell[c=2]{r} Total for Question 12 & & 6 & \\
\end{msscheme}

\msQuestion{13}
\begin{msscheme}{
  hline{3,5-7,9-12} = {2-Z}{0.4pt, msGrid},
  hline{4,8,13} = {0.4pt, msGrid},
  row{13} = {font=\bfseries},
}
Question & Scheme & Marks & Notes \\
13(a) & {Integrand undefined (unbounded) at the lower limit \(x=0\)} & E1 & {States the integrand is undefined, or tends to infinity, at \(x=0\). Not for a reason based on the range or the upper limit.} \\
\nopagebreak & & [1] & \\
13(b) & {\(\int x^{-\frac12}(\ln x)^2\,\mathrm{d}x = 2x^{\frac12}(\ln x)^2\)\\[4pt] \(\phantom{A} - \int 2x^{\frac12}\times\frac{2\ln x}{x}\,\mathrm{d}x\)} & M1 A1 & {\msNote{M1}{Integrates by parts with \(u=(\ln x)^2\), reaching the form \(px^{\frac12}(\ln x)^2 - q\int x^{-\frac12}\ln x\,\mathrm{d}x\). Condone coefficient slips.}\\ \msNote{A1}{Correct; need not be simplified.}} \\
& {\(\int x^{-\frac12}\ln x\,\mathrm{d}x = 2x^{\frac12}\ln x - \int 2x^{-\frac12}\,\mathrm{d}x\)} & dM1 & {Integrates \(\int x^{-\frac12}\ln x\,\mathrm{d}x\) by parts with \(u=\ln x\), reaching the form \(rx^{\frac12}\ln x - s\int x^{-\frac12}\,\mathrm{d}x\). Condone coefficient slips.} \\
& {\(= 2\sqrt{x}(\ln x)^2 - 4\left(2\sqrt{x}\ln x - 4\sqrt{x}\right) + c\)\\[4pt] \(= 2\sqrt{x}\left((\ln x)^2 - 4\ln x + 8\right) + c\)} & A1~AG & {Both stages correctly combined, final integral done, leading without error to the printed result, including \(+c\).} \\
\nopagebreak & & [4] & \\
13(c) & {\(\lim_{t\to0^{+}}\int_t^{e} x^{-\frac12}(\ln x)^2\,\mathrm{d}x\)\\[4pt] \(\phantom{A} = \lim_{t\to0^{+}}\left[2\sqrt{x}\left((\ln x)^2-4\ln x+8\right)\right]_t^{e}\)} & M1 & {Replaces the lower limit by \(t\), or another letter, and substitutes both limits into the result of (b). Not for substituting \(x=0\).} \\
& {\(= 10\sqrt{e} - 2\sqrt{t}\left((\ln t)^2 - 4\ln t + 8\right)\)} & A1 & {Correct expression in \(t\), oe, e.g. expanded.} \\
& {As \(t\to0^{+}\): \(\sqrt{t}(\ln t)^2\to0\), \(\sqrt{t}\ln t\to0\), \(\sqrt{t}\to0\)} & dM1 & {Lets \(t\to0^{+}\), using the given result for the \(\sqrt{t}(\ln t)^2\) and \(\sqrt{t}\ln t\) terms and \(\sqrt{t}\to0\).} \\
& {Limit finite, so integral exists, \(= 10\sqrt{e}\)} & A1 & {cso. Exact \(10\sqrt{e}\) with a statement that the limit is finite, or the integral exists.} \\
\nopagebreak & & [4] & \\
\nopagebreak \SetCell[c=2]{r} Total for Question 13 & & 9 & \\
\end{msscheme}

\msQuestion{14}
\begin{msscheme}{
  hline{3-7,9,11-12} = {2-Z}{0.4pt, msGrid},
  hline{8,10,13} = {0.4pt, msGrid},
  row{13} = {font=\bfseries},
}
Question & Scheme & Marks & Notes \\
14(a) & {\(\int \frac{(\ln x)^2}{x^2}\,\mathrm{d}x = -\frac{(\ln x)^2}{x} + \int \frac{2\ln x}{x^2}\,\mathrm{d}x\)} & M1 A1 & {\msNote{M1}{Integrates by parts with \(u=(\ln x)^2\) (or \(t^2\) after \(x=e^t\)) to the form \(p\frac{(\ln x)^2}{x} + q\int\frac{\ln x}{x^2}\,\mathrm{d}x\). Condone sign and coefficient slips.}\\ \msNote{A1}{Correct, need not be simplified.}} \\
& {\(\int \frac{\ln x}{x^2}\,\mathrm{d}x = -\frac{\ln x}{x} + \int \frac{1}{x^2}\,\mathrm{d}x\)} & dM1 & {Integrates \(\frac{\ln x}{x^2}\) (or \(te^{-t}\)) by parts to the form \(p\frac{\ln x}{x} + q\int\frac{1}{x^2}\,\mathrm{d}x\). Condone sign and coefficient slips.} \\
& {\(-\frac{(\ln x)^2}{x} - \frac{2\ln x}{x} - \frac{2}{x}\)} & A1 & {oe; constant not needed.} \\
& {\(\left(-\frac{(\ln a)^2}{a} - \frac{2\ln a}{a} - \frac{2}{a}\right) - (-2)\)} & dM1 & {Substitutes \(a\) and 1 into their antiderivative from two integrations by parts and subtracts.} \\
& {\(2 - \frac{(\ln a)^2}{a} - \frac{2\ln a}{a} - \frac{2}{a}\)} & A1 & {oe} \\
\nopagebreak & & [6] & \\
14(b) & {\(0\)} & B1 & \\
\nopagebreak & & [1] & \\
14(c) & {\(\lim_{a\to\infty}\left(2 - \frac{(\ln a)^2}{a} - \frac{2\ln a}{a} - \frac{2}{a}\right)\)\\[4pt] \(\phantom{A} = 2 - 0 - 0 - 0\)} & M1 & {Lets \(a\to\infty\) in their (a) of the form constant plus multiples of \(\frac{(\ln a)^2}{a}\), \(\frac{\ln a}{a}\), \(\frac1a\), stating each such term tends to 0.} \\
& {\(\int_1^{\infty} \frac{(\ln x)^2}{x^2}\,\mathrm{d}x = 2\)} & A1 & \\
\nopagebreak & & [2] & \\
\nopagebreak \SetCell[c=2]{r} Total for Question 14 & & 9 & \\
\end{msscheme}

\msQuestion{15}
\begin{msscheme}{
  hline{3-9} = {2-Z}{0.4pt, msGrid},
  hline{10} = {0.4pt, msGrid},
  row{10} = {font=\bfseries},
}
Question & Scheme & Marks & Notes \\
15 & {\(\dfrac{5x-3}{(x^2+3)(x+1)}=\dfrac{A}{x+1}+\dfrac{Bx+C}{x^2+3}\)} & M1 & {Uses a constant over \(x+1\) and a linear numerator over \(x^2+3\).} \\
& {\(5x-3=A(x^2+3)+(Bx+C)(x+1)\)\\[4pt] \(\Rightarrow A=\ldots,\ B=\ldots,\ C=\ldots\)} & dM1 & {Multiplies through and finds all three constants by substitution or equating coefficients. Condone arithmetic slips.} \\
& {\(A=-2,\ B=2,\ C=3\)} & A1 & \\
& {\(\int\left(-\dfrac{2}{x+1}+\dfrac{2x}{x^2+3}\right)\mathrm{d}x\)\\[4pt] \(\phantom{=}=-2\ln(x+1)+\ln(x^2+3)\)} & M1 & {Integrates their \(\frac{A}{x+1}\) and \(\frac{Bx}{x^2+3}\) to the forms \(\alpha\ln(x+1)\) and \(\beta\ln(x^2+3)\). Condone coefficient slips.} \\
& {\(\int\dfrac{3}{x^2+3}\,\mathrm{d}x=\dfrac{3}{\sqrt{3}}\arctan\left(\dfrac{x}{\sqrt{3}}\right)\)} & M1 & {Integrates their \(\frac{C}{x^2+3}\) to the form \(k\arctan\left(\frac{x}{\sqrt{3}}\right)\).} \\
& {\(\ln\left(\dfrac{x^2+3}{(x+1)^2}\right)+\sqrt{3}\arctan\left(\dfrac{x}{\sqrt{3}}\right)\)} & A1 & {Correct antiderivative oe, logarithms combined or not; modulus signs condoned.} \\
& {\(\lim_{t\to\infty}\left[\ln\left(\dfrac{t^2+3}{(t+1)^2}\right)+\sqrt{3}\arctan\left(\dfrac{t}{\sqrt{3}}\right)\right]\)\\[4pt] \(\phantom{=}-(\ln 3+0)\)\\[4pt] \(=0+\sqrt{3}\cdot\dfrac{\pi}{2}-\ln 3\)} & dM1 & {Replaces \(\infty\) by a limit, applies both limits, combines the logarithms into one tending to \(\ln 1\), and uses \(\arctan\to\frac{\pi}{2}\). Not for evaluating separate logarithms at infinity.} \\
& {\(\tfrac12\sqrt{3}\,\pi+\ln\tfrac13\): \(p=\tfrac12,\ q=\tfrac13\)} & A1 & {cso. Needs \(q=\frac13\) or \(\ln\frac13\); \(-\ln 3\) alone scores A0.} \\
\nopagebreak \SetCell[c=2]{r} Total for Question 15 & & 8 & \\
\end{msscheme}

\msQuestion{16}
\begin{msscheme}{
  hline{3,5-7,9-12,14} = {2-Z}{0.4pt, msGrid},
  hline{4,8,13,15} = {0.4pt, msGrid},
  row{15} = {font=\bfseries},
}
Question & Scheme & Marks & Notes \\
16(a) & {As \(y \to 0^{+}\), \(\ln y \to -\infty\), \((\ln y)^2 \to \infty\)\\[4pt] Integrand unbounded at \(y = 0\)} & E1 & {States the integrand is unbounded or undefined at \(y = 0\), referring to \(\ln y\) or \((\ln y)^2\) there.} \\
\nopagebreak & & [1] & \\
16(b)(i) & {\((\ln y)^2 = x^2\), \(\mathrm{d}y = -e^{-x}\,\mathrm{d}x\)} & M1 & {Replaces \((\ln y)^2\) by \(x^2\) and \(\mathrm{d}y\) by \(\pm e^{-x}\,\mathrm{d}x\). Condone a sign slip in \(\mathrm{d}y\).} \\
& {\(y = 1 \Rightarrow x = 0\)\\[4pt] \(y \to 0^{+} \Rightarrow x = -\ln y \to \infty\)} & B1 & {Both limits, with \(x \to \infty\) justified as \(y \to 0^{+}\) by \(-\ln y \to \infty\) or \(e^{-x} \to 0\). B0 for \(y = 0 \Rightarrow x = \infty\) with no limiting argument.} \\
& {\(\int_{\infty}^{0} x^2\left(-e^{-x}\right)\mathrm{d}x = \int_{0}^{\infty} x^2 e^{-x}\,\mathrm{d}x\)} & A1~AG & {Correct substituted integral with limits \(\infty\) to \(0\) and the minus sign, reversed to the printed result without error.} \\
\nopagebreak & & [3] & \\
16(b)(ii) & {\(\int x^2e^{-x}\,\mathrm{d}x = -x^2e^{-x} + \int 2xe^{-x}\,\mathrm{d}x\)} & M1 A1 & {\msNote{M1}{Integrates by parts with \(u = x^2\), reaching \(\pm x^2e^{-x} \pm \int kxe^{-x}\,\mathrm{d}x\). Tabular method accepted.}} \\
& {\(\phantom{\int} = -x^2e^{-x} - 2xe^{-x} - 2e^{-x}\)} & dM1 A1 & {\msNote{dM1}{Integrates their \(\int kxe^{-x}\,\mathrm{d}x\) by parts, reaching the form \(pxe^{-x} + qe^{-x}\).}\\ \msNote{A1}{oe; constant not needed.}} \\
& {\(\left[-e^{-x}(x^2+2x+2)\right]_0^t = 2 - e^{-t}(t^2+2t+2)\)\\[4pt] As \(t \to \infty\), \(t^2e^{-t}, te^{-t}, e^{-t} \to 0\)} & dM1 & {Applies limits \(0\) and \(t\) to their antiderivative of the form \(e^{-x}\times\) quadratic, and lets \(t \to \infty\) with the exponential terms tending to 0. Not for substituting \(\infty\) directly.} \\
& {\(\int_{0}^{\infty} x^2 e^{-x}\,\mathrm{d}x = 2\)} & A1 & {cso. A bare answer scores nothing.} \\
\nopagebreak & & [6] & \\
16(b)(iii) & {\(\int_{0}^{1} (\ln y)^2\,\mathrm{d}y = 2\)} & B1ft & {ft their value from (b)(ii).} \\
\nopagebreak & & [1] & \\
\nopagebreak \SetCell[c=2]{r} Total for Question 16 & & 11 & \\
\end{msscheme}

\end{document}
